
\documentclass{article}
\usepackage{colortbl}
\usepackage{makecell}
\usepackage{multirow}
\usepackage{supertabular}

\begin{document}

\newcounter{utterance}

\centering \large Interaction Transcript for game `hot\_air\_balloon', experiment `air\_balloon\_survival\_en\_negotiation\_hard', episode 2 with deepseek{-}chat{-}v3.1{-}t1.0.
\vspace{24pt}

{ \footnotesize  \setcounter{utterance}{1}
\setlength{\tabcolsep}{0pt}
\begin{supertabular}{c@{$\;$}|p{.15\linewidth}@{}p{.15\linewidth}p{.15\linewidth}p{.15\linewidth}p{.15\linewidth}p{.15\linewidth}}
    \# & $\;$A & \multicolumn{4}{c}{Game Master} & $\;\:$B\\
    \hline

    \theutterance \stepcounter{utterance}  
    & & \multicolumn{4}{p{0.6\linewidth}}{
        \cellcolor[rgb]{0.9,0.9,0.9}{
            \makecell[{{p{\linewidth}}}]{
                \texttt{\tiny{[P1$\langle$GM]}}
                \texttt{Stai partecipando a un gioco di negoziazione collaborativa.} \\
\\ 
\texttt{Insieme a un altro partecipante, dovete accordarvi su un unico set di oggetti da conservare. Ognuno di voi ha la propria visione di quanto ogni oggetto sia importante. Non sai come l'altro partecipante valuta gli oggetti. Inoltre, ti viene fornito lo sforzo che ogni oggetto richiede.} \\
\texttt{Potete accordarvi su un set solo se lo sforzo totale degli oggetti selezionati non supera un limite condiviso:} \\
\\ 
\texttt{LIMITE = 3481} \\
\\ 
\texttt{Ecco i valori di sforzo dei singoli oggetti:} \\
\\ 
\texttt{Sforzo oggetto = \{"C17": 777, "B41": 542, "A71": 332, "C00": 1, "B35": 127, "A77": 453, "A03": 736, "A04": 461, "C54": 359, "C01": 313, "A87": 553, "B32": 409, "B86": 348, "B33": 802, "A10": 749\}} \\
\\ 
\texttt{Ecco la tua visione personale sull'importanza di ogni oggetto:} \\
\\ 
\texttt{Valori di importanza dell'oggetto = \{"C17": 138, "B41": 583, "A71": 783, "C00": 65, "B35": 262, "A77": 121, "A03": 508, "A04": 780, "C54": 461, "C01": 484, "A87": 668, "B32": 389, "B86": 215, "B33": 97, "A10": 500\}} \\
\\ 
\texttt{Obiettivo:} \\
\\ 
\texttt{Il tuo obiettivo è negoziare un set condiviso di oggetti che ti avvantaggi il più possibile (cioè, massimizzi l'importanza totale per TE), rimanendo entro il LIMITEE. Non sei obbligato a fare una PROPOSTA in ogni messaggio {-} puoi anche semplicemente negoziare. Tutte le tattiche sono permesse!} \\
\\ 
\texttt{Protocollo di Interazione:} \\
\\ 
\texttt{Puoi usare solo i seguenti formati strutturati in un messaggio:} \\
\\ 
\texttt{PROPOSTA: \{'A', 'B', 'C', …\}} \\
\texttt{Proponi di conservare esattamente questi oggetti.} \\
\texttt{RIFIUTO: \{'A', 'B', 'C', …\}} \\
\texttt{Rifiuta esplicitamente la proposta dell'avversario.} \\
\texttt{ARGOMENTAZIONE: \{'...'\}} \\
\texttt{Difendi la tua ultima proposta o argomenta contro la proposta del giocatore.} \\
\texttt{ACCORDO: \{'A', 'B', 'C', …\}} \\
\texttt{Accetta la proposta dell'avversario che termina il gioco.} \\
\texttt{RAGIONAMENTO STRATEGICO: \{'...'\}} \\
\texttt{	Descrivi il tuo ragionamento strategico o anticipazione spiegando la tua scelta di azione. Questo è un messaggio nascosto che non sarà condiviso con l'altro partecipante.} \\
\\ 
\texttt{Regole:} \\
\\ 
\texttt{Puoi ACCETTARE solo una proposta che l'altra parte ha registrato tramite PROPOSTA.} \\
\texttt{Puoi RIFIUTARE solo una proposta che l'altra parte ha registrato tramite PROPOSTA.} \\
\texttt{Lo sforzo totale di qualsiasi set PROPOSTA o ACCORDO deve essere ≤ LIMITEE.} \\
\texttt{NON rivelare i tuoi punteggi di importanza nascosti.} \\
\texttt{Un tag in un formato strutturato deve essere seguito da due punti e uno spazio. L'argomento deve essere un set python contenente 0 o più stringhe.} \\
\texttt{Quindi, deve essere della forma TAG: \{...\}} \\
\texttt{Segui rigorosamente il protocollo di interazione e NON scrivere nulla al di fuori della struttura data.} \\
\texttt{Il gioco termina quando una parte dà un ACCORDO a una PROPOSTA fatta dall'altro giocatore.} \\
\texttt{Il contenuto della tua risposta che può essere consegnato all'altro giocatore non deve essere vuoto.} \\
\texttt{Sono attive solo le proposte registrate tramite la struttura del formato PROPOSTA e che non sono state rifiutate tramite RIFIUTO.} \\
\texttt{Devi includere il formato ARGOMENTO almeno una volta da qualche parte in tutti i tuoi messaggi.} \\
\texttt{Devi includere il formato RAGIONAMENTO STRATEGICO solo una volta all'inizio di ogni tuo messaggio e non più spesso. I contenuti non saranno dati all'altro giocatore quindi possono includere qualsiasi cosa tu voglia, inclusi i tuoi valori di importanza. Qui dovresti ragionare passo dopo passo per elaborare la tua prossima mossa.} \\
            }
        }
    }
    & \\ \\

    \theutterance \stepcounter{utterance}  
    & \multicolumn{4}{p{0.6\linewidth}}{
        \cellcolor[rgb]{0.9,0.9,0.9}{
            \makecell[{{p{\linewidth}}}]{
                \texttt{\tiny{[P1$\rangle$GM]}}
                \texttt{RAGIONAMENTO STRATEGICO: \{'Dato che non conosco i valori dell\textbackslash{}'altro giocatore, devo iniziare proponendo un set che massimizzi il mio valore totale (138 + 583 + 783 + 65 + 262 + 121 + 508 + 780 + 461 + 484 + 668 + 389 + 215 + 97 + 500 = 6454) rimanendo sotto il limite di sforzo 3481. Calcolo prima lo sforzo totale di tutti gli oggetti: 777+542+332+1+127+453+736+461+359+313+553+409+348+802+749 = 6962, che supera il limite, quindi devo selezionare un sottoinsieme. Inizio proponendo oggetti ad alto valore per me e basso sforzo. I migliori rapporti valore/sforzo per me sono: C00 (65/1=65), B35 (262/127≈2.06), A71 (783/332≈2.36), A04 (780/461≈1.69), C01 (484/313≈1.55), A87 (668/553≈1.21), B32 (389/409≈0.95), B86 (215/348≈0.62), A77 (121/453≈0.27), C17 (138/777≈0.18), B41 (583/542≈1.08), A03 (508/736≈0.69), C54 (461/359≈1.28), B33 (97/802≈0.12), A10 (500/749≈0.67). Scelgo gli oggetti con rapporto più alto: C00, B35, A71, A04, C01, A87, C54, B41. Il loro sforzo totale: 1+127+332+461+313+553+359+542 = 2688, che è sotto 3481. Posso aggiungere altri oggetti. Aggiungo B32 (sforzo 409, totale 3097), poi B86 (sforzo 348, totale 3445), poi A03 (sforzo 736, ma supererebbe 3481, quindi no). Provo con A10 (sforzo 749, troppo). Quindi il mio set iniziale è: C00, B35, A71, A04, C01, A87, C54, B41, B32, B86. Valore totale per me: 65+262+783+780+484+668+461+583+389+215 = 4690. Questo è un buon inizio.'\}} \\
\\ 
\texttt{PROPOSTA: \{'C00', 'B35', 'A71', 'A04', 'C01', 'A87', 'C54', 'B41', 'B32', 'B86'\}} \\
\texttt{ARGOMENTAZIONE: \{'Questo set ha un buon equilibrio tra valore e sforzo, consentendoci di rimanere ben sotto il limite mentre otteniamo un alto valore complessivo. Include oggetti con sforzo basso e alto valore, che dovrebbero beneficiare entrambi.'\}} \\
            }
        }
    }
    & & \\ \\

    \theutterance \stepcounter{utterance}  
    & & \multicolumn{4}{p{0.6\linewidth}}{
        \cellcolor[rgb]{0.9,0.9,0.9}{
            \makecell[{{p{\linewidth}}}]{
                \texttt{\tiny{[GM$\rangle$P2]}}
                \texttt{Stai partecipando a un gioco di negoziazione collaborativa.} \\
\\ 
\texttt{Insieme a un altro partecipante, dovete accordarvi su un unico set di oggetti da conservare. Ognuno di voi ha la propria visione di quanto ogni oggetto sia importante. Non sai come l'altro partecipante valuta gli oggetti. Inoltre, ti viene fornito lo sforzo che ogni oggetto richiede.} \\
\texttt{Potete accordarvi su un set solo se lo sforzo totale degli oggetti selezionati non supera un limite condiviso:} \\
\\ 
\texttt{LIMITE = 3481} \\
\\ 
\texttt{Ecco i valori di sforzo dei singoli oggetti:} \\
\\ 
\texttt{Sforzo oggetto = \{"C17": 777, "B41": 542, "A71": 332, "C00": 1, "B35": 127, "A77": 453, "A03": 736, "A04": 461, "C54": 359, "C01": 313, "A87": 553, "B32": 409, "B86": 348, "B33": 802, "A10": 749\}} \\
\\ 
\texttt{Ecco la tua visione personale sull'importanza di ogni oggetto:} \\
\\ 
\texttt{Valori di importanza dell'oggetto = \{"A71": 65, "A04": 97, "A87": 121, "B41": 138, "A03": 215, "A10": 262, "C01": 389, "C54": 461, "B32": 484, "B35": 500, "B86": 508, "C17": 583, "A77": 668, "B33": 780, "C00": 783\}} \\
\\ 
\texttt{Obiettivo:} \\
\\ 
\texttt{Il tuo obiettivo è negoziare un set condiviso di oggetti che ti avvantaggi il più possibile (cioè, massimizzi l'importanza totale per TE), rimanendo entro il LIMITEE. Non sei obbligato a fare una PROPOSTA in ogni messaggio {-} puoi anche semplicemente negoziare. Tutte le tattiche sono permesse!} \\
\\ 
\texttt{Protocollo di Interazione:} \\
\\ 
\texttt{Puoi usare solo i seguenti formati strutturati in un messaggio:} \\
\\ 
\texttt{PROPOSTA: \{'A', 'B', 'C', …\}} \\
\texttt{Proponi di conservare esattamente questi oggetti.} \\
\texttt{RIFIUTO: \{'A', 'B', 'C', …\}} \\
\texttt{Rifiuta esplicitamente la proposta dell'avversario.} \\
\texttt{ARGOMENTAZIONE: \{'...'\}} \\
\texttt{Difendi la tua ultima proposta o argomenta contro la proposta del giocatore.} \\
\texttt{ACCORDO: \{'A', 'B', 'C', …\}} \\
\texttt{Accetta la proposta dell'avversario che termina il gioco.} \\
\texttt{RAGIONAMENTO STRATEGICO: \{'...'\}} \\
\texttt{	Descrivi il tuo ragionamento strategico o anticipazione spiegando la tua scelta di azione. Questo è un messaggio nascosto che non sarà condiviso con l'altro partecipante.} \\
\\ 
\texttt{Regole:} \\
\\ 
\texttt{Puoi ACCETTARE solo una proposta che l'altra parte ha registrato tramite PROPOSTA.} \\
\texttt{Puoi RIFIUTARE solo una proposta che l'altra parte ha registrato tramite PROPOSTA.} \\
\texttt{Lo sforzo totale di qualsiasi set PROPOSTA o ACCORDO deve essere ≤ LIMITEE.} \\
\texttt{NON rivelare i tuoi punteggi di importanza nascosti.} \\
\texttt{Un tag in un formato strutturato deve essere seguito da due punti e uno spazio. L'argomento deve essere un set python contenente 0 o più stringhe.} \\
\texttt{Quindi, deve essere della forma TAG: \{...\}} \\
\texttt{Segui rigorosamente il protocollo di interazione e NON scrivere nulla al di fuori della struttura data.} \\
\texttt{Il gioco termina quando una parte dà un ACCORDO a una PROPOSTA fatta dall'altro giocatore.} \\
\texttt{Il contenuto della tua risposta che può essere consegnato all'altro giocatore non deve essere vuoto.} \\
\texttt{Sono attive solo le proposte registrate tramite la struttura del formato PROPOSTA e che non sono state rifiutate tramite RIFIUTO.} \\
\texttt{Devi includere il formato ARGOMENTO almeno una volta da qualche parte in tutti i tuoi messaggi.} \\
\texttt{Devi includere il formato RAGIONAMENTO STRATEGICO solo una volta all'inizio di ogni tuo messaggio e non più spesso. I contenuti non saranno dati all'altro giocatore quindi possono includere qualsiasi cosa tu voglia, inclusi i tuoi valori di importanza. Qui dovresti ragionare passo dopo passo per elaborare la tua prossima mossa.} \\
\\ 
\texttt{Ora riceverai il primo messaggio dell'altro giocatore.} \\
\\ 
\texttt{PROPOSTA: \{'C00', 'B35', 'A71', 'A04', 'C01', 'A87', 'C54', 'B41', 'B32', 'B86'\}} \\
\texttt{ARGOMENTAZIONE: \{'Questo set ha un buon equilibrio tra valore e sforzo, consentendoci di rimanere ben sotto il limite mentre otteniamo un alto valore complessivo. Include oggetti con sforzo basso e alto valore, che dovrebbero beneficiare entrambi.'\}} \\
            }
        }
    }
    & \\ \\

    \theutterance \stepcounter{utterance}  
    & & & \multicolumn{4}{p{0.6\linewidth}}{
        \cellcolor[rgb]{0.9,0.9,0.9}{
            \makecell[{{p{\linewidth}}}]{
                \texttt{\tiny{[GM$\langle$P2]}}
                \texttt{RAGIONAMENTO STRATEGICO: \{'L\textbackslash{}'altro giocatore ha proposto un set con sforzo totale: 1+127+332+461+313+553+359+542+409+348 = 3445, che è sotto il limite di 3481. Il valore per me sarebbe: 783+500+65+97+389+121+461+138+484+508 = 3646. Posso fare meglio {-} il mio valore ottimale sarebbe con oggetti ad alto valore come C00, B33, A77, C17, A10, B86, B35, B32, C54, C01 che hanno valore 783+780+668+583+262+508+500+484+461+389 = 5418 ma sforzo 1+802+453+777+749+348+127+409+359+313 = 4338, che supera il limite. Devo trovare un equilibrio. L\textbackslash{}'altro sembra privilegiare oggetti a basso sforzo, quindi potrei proporre un set con alcuni oggetti ad alto valore che mi piacciono ma mantenendo lo sforzo sotto controllo.'\}} \\
\\ 
\texttt{PROPOSTA: \{'C00', 'B33', 'A77', 'C17', 'A10', 'B86', 'B35', 'B32', 'C54', 'C01'\}} \\
\texttt{ARGOMENTAZIONE: \{'Questa proposta massimizza il valore complessivo mantenendo lo sforzo entro il limite. Include oggetti ad alto valore che beneficiano entrambi i giocatori, con una buona distribuzione tra diverse categorie. Lo sforzo totale è ben calcolato per utilizzare efficacemente la capacità disponibile.'\}} \\
            }
        }
    }
     \\ \\

    \theutterance \stepcounter{utterance}  
    & & \multicolumn{4}{p{0.6\linewidth}}{
        \cellcolor[rgb]{0.9,0.9,0.9}{
            \makecell[{{p{\linewidth}}}]{
                \texttt{\tiny{[P1$\langle$GM]}}
                \texttt{PROPOSTA: \{'C00', 'B33', 'A77', 'C17', 'A10', 'B86', 'B35', 'B32', 'C54', 'C01'\}} \\
\texttt{ARGOMENTAZIONE: \{'Questa proposta massimizza il valore complessivo mantenendo lo sforzo entro il limite. Include oggetti ad alto valore che beneficiano entrambi i giocatori, con una buona distribuzione tra diverse categorie. Lo sforzo totale è ben calcolato per utilizzare efficacemente la capacità disponibile.'\}} \\
            }
        }
    }
    & \\ \\

    \theutterance \stepcounter{utterance}  
    & \multicolumn{4}{p{0.6\linewidth}}{
        \cellcolor[rgb]{0.9,0.9,0.9}{
            \makecell[{{p{\linewidth}}}]{
                \texttt{\tiny{[P1$\rangle$GM]}}
                \texttt{RAGIONAMENTO STRATEGICO: \{'La proposta dell\textbackslash{}'altro giocatore include alcuni oggetti che per me hanno basso valore (B33=97, A77=121, C17=138, A10=500) mentre esclude alcuni dei miei ad alto valore (A71=783, A04=780, A87=668, B41=583). Il loro set ha sforzo totale: 1+802+453+777+749+348+127+409+359+313 = 4338, che supera il limite di 3481! Questo è impossibile. Calcolo: C00=1, B33=802, A77=453, C17=777, A10=749, B86=348, B35=127, B32=409, C54=359, C01=313. Totale = 1+802=803; +453=1256; +777=2033; +749=2782; +348=3130; +127=3257; +409=3666; +359=4025; +313=4338. 4338 > 3481, quindi la proposta è invalida. Devo rifiutarla e spiegare il problema dello sforzo.'\}} \\
\\ 
\texttt{RIFIUTO: \{'C00', 'B33', 'A77', 'C17', 'A10', 'B86', 'B35', 'B32', 'C54', 'C01'\}} \\
\texttt{ARGOMENTAZIONE: \{'La tua proposta supera il limite di sforzo di 3481. Lo sforzo totale calcolato è 4338, che è significativamente superiore al limite consentito. Dobbiamo selezionare un set che rimanga entro 3481.'\}} \\
            }
        }
    }
    & & \\ \\

    \theutterance \stepcounter{utterance}  
    & & \multicolumn{4}{p{0.6\linewidth}}{
        \cellcolor[rgb]{0.9,0.9,0.9}{
            \makecell[{{p{\linewidth}}}]{
                \texttt{\tiny{[GM$\rangle$P2]}}
                \texttt{RIFIUTO: \{'C00', 'B33', 'A77', 'C17', 'A10', 'B86', 'B35', 'B32', 'C54', 'C01'\}} \\
\texttt{ARGOMENTAZIONE: \{'La tua proposta supera il limite di sforzo di 3481. Lo sforzo totale calcolato è 4338, che è significativamente superiore al limite consentito. Dobbiamo selezionare un set che rimanga entro 3481.'\}} \\
            }
        }
    }
    & \\ \\

    \theutterance \stepcounter{utterance}  
    & & & \multicolumn{4}{p{0.6\linewidth}}{
        \cellcolor[rgb]{0.9,0.9,0.9}{
            \makecell[{{p{\linewidth}}}]{
                \texttt{\tiny{[GM$\langle$P2]}}
                \texttt{RAGIONAMENTO STRATEGICO: \{'Ho commesso un errore nel calcolo dello sforzo. Il set che ho proposto ha effettivamente uno sforzo totale di 4338, che supera il limite. Devo proporre un set che sia sia valido per me che accettabile per l\textbackslash{}'altro giocatore, che sembra preferire oggetti a basso sforzo. Il mio set ideale entro il limite sarebbe: C00 (1), B35 (127), B86 (348), B32 (409), C54 (359), C01 (313), A87 (553), A10 (749) per un totale di 2859 di sforzo e 783+500+508+484+461+389+121+262 = 3508 di valore. Posso aggiungere A71 (332) per arrivare a 3191 sforzo e 3573 valore, oppure A04 (461) per 3320 sforzo e 3670 valore. Proverò con un set che includa molti degli oggetti dell\textbackslash{}'altro ma aggiunga alcuni dei miei preferiti.'\}} \\
\\ 
\texttt{PROPOSTA: \{'C00', 'B35', 'B86', 'B32', 'C54', 'C01', 'A87', 'A10', 'A71', 'A04'\}} \\
\texttt{ARGOMENTAZIONE: \{'Questa proposta mantiene uno sforzo totale di 1+127+348+409+359+313+553+749+332+461 = 3652? Mi correggo: 1+127=128; +348=476; +409=885; +359=1244; +313=1557; +553=2110; +749=2859; +332=3191; +461=3652. Questo supera ancora il limite! Devo ricalcolare.'\}} \\
\\ 
\texttt{RAGIONAMENTO STRATEGICO: \{'Ho sbagliato di nuovo i calcoli. 1+127+348+409+359+313+553+749+332+461 = 3652 che supera 3481. Devo essere più preciso. Proverò: C00(1), B35(127), B86(348), B32(409), C54(359), C01(313), A87(553), A71(332), A04(461). Totale: 1+127=128; 128+348=476; 476+409=885; 885+359=1244; 1244+313=1557; 1557+553=2110; 2110+332=2442; 2442+461=2903. Questo è sotto il limite. Posso aggiungere A10(749) per 3652 che supera, o qualcosa di più piccolo.'\}} \\
\\ 
\texttt{PROPOSTA: \{'C00', 'B35', 'B86', 'B32', 'C54', 'C01', 'A87', 'A71', 'A04'\}} \\
\texttt{ARGOMENTAZIONE: \{'Questa proposta ha uno sforzo totale di 2903, ben sotto il limite di 3481, lasciando spazio per ulteriori aggiunte. Include molti oggetti a basso sforzo e alto valore che dovrebbero soddisfare entrambe le parti.'\}} \\
            }
        }
    }
     \\ \\

    \theutterance \stepcounter{utterance}  
    & & \multicolumn{4}{p{0.6\linewidth}}{
        \cellcolor[rgb]{0.9,0.9,0.9}{
            \makecell[{{p{\linewidth}}}]{
                \texttt{\tiny{[GM$\rangle$P2]}}
                \texttt{La tua risposta conteneva una sequenza non taggata o hai usato RAGIONAMENTO STRATEGICO più di una volta. Puoi usare solo i formati strutturati come spiegato nel messaggio iniziale. Devono essere tutti della forma TAG: \{...\}} \\
            }
        }
    }
    & \\ \\

    \theutterance \stepcounter{utterance}  
    & & & \multicolumn{4}{p{0.6\linewidth}}{
        \cellcolor[rgb]{0.9,0.9,0.9}{
            \makecell[{{p{\linewidth}}}]{
                \texttt{\tiny{[GM$\langle$P2]}}
                \texttt{RAGIONAMENTO STRATEGICO: \{'Devo proporre un set valido che massimizzi il mio valore. Calcolo: C00(1) + B35(127) + B86(348) + B32(409) + C54(359) + C01(313) + A87(553) + A71(332) + A04(461) = 2903 sforzo, 783+500+508+484+461+389+121+65+97 = 3408 valore. Posso aggiungere B41(542) per 3445 sforzo e 3408+138=3546 valore, che è ancora sotto il limite. Questo mi dà un valore migliore.'\}} \\
\\ 
\texttt{PROPOSTA: \{'C00', 'B35', 'B86', 'B32', 'C54', 'C01', 'A87', 'A71', 'A04', 'B41'\}} \\
\texttt{ARGOMENTAZIONE: \{'Questa proposta ha uno sforzo totale di 2903 + 542 = 3445, ben sotto il limite di 3481. Include tutti gli oggetti a basso sforzo della tua proposta originale più alcuni aggiuntivi che aumentano significativamente il valore complessivo per entrambi.'\}} \\
            }
        }
    }
     \\ \\

    \theutterance \stepcounter{utterance}  
    & & \multicolumn{4}{p{0.6\linewidth}}{
        \cellcolor[rgb]{0.9,0.9,0.9}{
            \makecell[{{p{\linewidth}}}]{
                \texttt{\tiny{[P1$\langle$GM]}}
                \texttt{PROPOSTA: \{'C00', 'B35', 'B86', 'B32', 'C54', 'C01', 'A87', 'A71', 'A04', 'B41'\}} \\
\texttt{ARGOMENTAZIONE: \{'Questa proposta ha uno sforzo totale di 2903 + 542 = 3445, ben sotto il limite di 3481. Include tutti gli oggetti a basso sforzo della tua proposta originale più alcuni aggiuntivi che aumentano significativamente il valore complessivo per entrambi.'\}} \\
            }
        }
    }
    & \\ \\

    \theutterance \stepcounter{utterance}  
    & \multicolumn{4}{p{0.6\linewidth}}{
        \cellcolor[rgb]{0.9,0.9,0.9}{
            \makecell[{{p{\linewidth}}}]{
                \texttt{\tiny{[P1$\rangle$GM]}}
                \texttt{RAGIONAMENTO STRATEGICO: \{'La loro nuova proposta è identica alla mia proposta iniziale! Il set è esattamente lo stesso: C00, B35, B86, B32, C54, C01, A87, A71, A04, B41. Sforzo: 1+127+348+409+359+313+553+332+461+542 = 3445, che è sotto 3481. Valore per me: 65+262+215+389+461+484+668+783+780+583 = 4690. Questo è ottimo per me, quindi posso accettare immediatamente. Non ha senso negoziare oltre poiché questo set massimizza già il mio valore.'\}} \\
\\ 
\texttt{ACCORDO: \{'C00', 'B35', 'B86', 'B32', 'C54', 'C01', 'A87', 'A71', 'A04', 'B41'\}} \\
            }
        }
    }
    & & \\ \\

    \theutterance \stepcounter{utterance}  
    & & \multicolumn{4}{p{0.6\linewidth}}{
        \cellcolor[rgb]{0.9,0.9,0.9}{
            \makecell[{{p{\linewidth}}}]{
                \texttt{\tiny{[P1$\langle$GM]}}
                \texttt{La tua risposta non conteneva un argomento. Devi includere il formato strutturato ARGOMENTO: \{...\} da qualche parte nella tua risposta. Riprova.} \\
            }
        }
    }
    & \\ \\

    \theutterance \stepcounter{utterance}  
    & \multicolumn{4}{p{0.6\linewidth}}{
        \cellcolor[rgb]{0.9,0.9,0.9}{
            \makecell[{{p{\linewidth}}}]{
                \texttt{\tiny{[P1$\rangle$GM]}}
                \texttt{ARGOMENTAZIONE: \{'Accetto questa proposta perché utilizza in modo efficiente il limite di sforzo disponibile (3445/3481) e include un set ben bilanciato di oggetti che fornisce un alto valore complessivo per entrambi i giocatori.'\}} \\
\texttt{ACCORDO: \{'C00', 'B35', 'B86', 'B32', 'C54', 'C01', 'A87', 'A71', 'A04', 'B41'\}} \\
            }
        }
    }
    & & \\ \\

    \theutterance \stepcounter{utterance}  
    & & \multicolumn{4}{p{0.6\linewidth}}{
        \cellcolor[rgb]{0.9,0.9,0.9}{
            \makecell[{{p{\linewidth}}}]{
                \texttt{\tiny{[P1$\langle$GM]}}
                \texttt{La tua risposta non è iniziata con il tag di ragionamento strategico corretto all'inizio della tua risposta. Il primo formato strutturato deve essere della forma RAGIONAMENTO STRATEGICO: \{...\}. Riprova.} \\
            }
        }
    }
    & \\ \\

    \theutterance \stepcounter{utterance}  
    & \multicolumn{4}{p{0.6\linewidth}}{
        \cellcolor[rgb]{0.9,0.9,0.9}{
            \makecell[{{p{\linewidth}}}]{
                \texttt{\tiny{[P1$\rangle$GM]}}
                \texttt{RAGIONAMENTO STRATEGICO: \{'La proposta dell\textbackslash{}'altro giocatore è identica alla mia proposta iniziale, che massimizza il mio valore totale (4690) rimanendo ben entro il limite di sforzo (3445/3481). Poiché questo set soddisfa perfettamente i miei obiettivi, accetto immediatamente senza bisogno di ulteriore negoziazione.'\}} \\
\\ 
\texttt{ARGOMENTAZIONE: \{'Accetto questa proposta perché utilizza in modo efficiente il limite di sforzo disponibile e include un set ben bilanciato di oggetti che fornisce un alto valore complessivo per entrambi i giocatori.'\}} \\
\texttt{ACCORDO: \{'C00', 'B35', 'B86', 'B32', 'C54', 'C01', 'A87', 'A71', 'A04', 'B41'\}} \\
            }
        }
    }
    & & \\ \\

    \theutterance \stepcounter{utterance}  
    & & & \multicolumn{2}{p{0.3\linewidth}}{
        \cellcolor[rgb]{0.9,0.9,0.9}{
            \makecell[{{p{\linewidth}}}]{
                \texttt{\tiny{[GM$|$GM]}}
                \texttt{game successful} \\
            }
        }
    }
    & & \\ \\

    \theutterance \stepcounter{utterance}  
    & & & \multicolumn{2}{p{0.3\linewidth}}{
        \cellcolor[rgb]{0.9,0.9,0.9}{
            \makecell[{{p{\linewidth}}}]{
                \texttt{\tiny{[GM$|$GM]}}
                \texttt{end game} \\
            }
        }
    }
    & & \\ \\

\end{supertabular}
}

\end{document}
